\begin{chapterbox}
    \chapter{Ein folgenschwerer Tod (2022)}
    \label{2022-ein-folgenschwerer-tod}
    \az{Jahr 65}

    \begin{center}
        Zwischenspiel der Magischen Abenteuer

        Fortsetzung von \hyperref[2022-der-verschwundene-feuertakuri]{Der verschwundene Feuertakuri (2022)}
    \end{center}
    
    Tarok, der letzte Drache, haucht sein Leben aus. Sein Tod hat bedeutende Konsequenzen weit über das Königreich Andor hinaus. Selbst im fernen Tulgor spüren Helden die Folgen. 

    \trennzeichen
    
    \nurvollversion{
        % https://legenden-von-andor.de/forum/viewtopic.php?p=94100#p94100
        \textit{eine sehr schöne Gelegenheit für jede Menge Aspielungen zu deinen diversen anderen Geschichten (und zum Ewigen Rat).}
        \newline
        – TroII
    }
\end{chapterbox}



\az{Jahr 65}


Ein Kopf, der den Unterirdischen Krieg mit eigenen, blutroten Augen miterlebt hatte, sank tief in den blutgetränken Boden Andors ein. Er schmeckte Matsch, Dreck und Blut. Ein letztes Ächzen drang aus dem riesigen Maul.

Damit schied Tarok, der letzte Drache, für immer aus dem Leben. Seine Flamme des Zorns erlosch.

Südlich davon, tief, tief unter den Bergen des Grauen Gebirges, wogte in einer gewaltigen Höhle ein Geflecht aus Düsternis umher. In der Mitte ragte ein riesiger Baum empor. Schimmernde Edelsteine ragten hier und da aus seiner längst versteinerten Oberfläche. Doch klar erkennbar unter dem durchscheinenden hellen Gestein floss in seinen dicken Adern stetig rotes Blut durch den Baum. Schon seit Jahrtausenden, seit der Erschaffung durch die uralten Drachen, hatte der Baum des Blutes hier in Krahal gestanden, den Drachen Leben gespendet und ihnen Kraft und Energie verliehen.

Nun stockte der Herzschlag des Baumes. Das rote Blut gerann. Ein Glimmen lief über die Rinde des Baums und verwandelte helles, lebendiges Gestein in tote, schwarze Asche. Der Baum verödete innert Minuten. Die Schreie tausender Drachenseelen erklangen ein letztes Mal und verstummten dann für immer.

Einen Augenblick lang, regte sich nichts mehr, als wäre die Zeit selbst in Krahal stehen geblieben.

Dann, urplötzlich, wurde Krahal zutiefst erschüttert. Das Geflecht der Düsternis wirbelte wild umher und zerriss die erstarrten Überreste des blutigen Baumes. Krahal verschlang sich selbst. Die gewaltige Höhle stürzte in sich zusammen.

Das Splittern von Stein hallte in den Bergtälern wider und übertönte die Musik der Toten in den Schluchten des Grauen Gebirges. Felsen barsten und Lawinen krachten. Erschütterungen gingen vom Epizentrum aus. Sie zogen sich in Wellen durchs Graue Gebirge, jagten durch Fels und Stein, zwischen Schluchten und durch die unterirdischen Gänge der Schildzwerge.

An der Ruine der Winterburg stürzte eine weitere Mauer unter lautem Gekrache in sich zusammen.

Rhona und Grone, die Stammesältesten der Agren, fielen in eine stille Umarmung, während ihre Höhle um sie herum bebte.

Schildzwerge fielen auf ihre Knie und sandten Stoßgebete aus, während ihre Umgebung zitterte und die Tiefminen Flammen spuckten.

Der Mechaniker Kjall raste in seinem kleinen Museum umher und versuchte unter Einsatz eines seiner Tiefminen-Golems verzweifelt, seine Schätze – Fackeln von Cavern, schöne Schilde, ein zersplittertes Hadrisches Stundenglas – auf ihren angestammten Sockeln zu behalten.

Der Riesenkrake Irlok zuckte unruhig umher, während gewaltige Wellen ans Ufer des Geheimen Sees brandeten und Stalaktiten von der Decke ins Wasser klatschten.

Die von Krahal ausgehenden Beben wogten auch bis ins Kuolema-Gebirge.

Tulgorische Höhlenwichte sprangen beiseite, als ihre Gänge sich teilten.

Ein tiefer Spalt tat sich im ewigen Eis auf.

Siantari, die oberste Dämonin des ewigen Eises, breitete ihre Arme aus und rief eine wirbelnde Wolke aus Schnee und Eis zu sich. Der Wirbel hob sie in die Höhe und ließ sie einige Meter über der zitternden Eisfläche schweben, weg von den Vibrationen der Erde. Ihre eisblauen Augen starrten regungslos auf das Felsentor, das ihr nun schon seit 49 Dekaden den Zugang zur restlichen Welt verwehrte.

Das Felsentor erzitterte wie der gesamte Rest des Kuolema-Gebirges. Staub rieselte herab und Splitter bröckelten. Doch die Magie der Temm, die das stabile Tor vor Jahrtausenden vor dem ersten Dämon des ewigen Eises verschlossen hatten, ließ nicht nach.

Dann war es vorbei. Die Welt, das Gebirge, das ewige Eis, sie alle standen wieder ruhig da. Krahal war nicht mehr. Shan, die Schattenhexe, blickte verwirrt um sich, als sie die gewaltige Seelenkraft Krahals tief unter der Erde auf einmal nicht mehr fühlen konnte.

Doch das Felsentor vor dem ewigen Eis des Kuolema hatte standgehalten. Die Eisdämonen waren immer noch in der Schlucht dahinter gefangen.

„Fećht!“, fluchte Siantari. So schnell, wie der Ärger in ihr aufgeschwellt war, wurde er auch schon wieder von der allgegenwärtigen Ruhe des Eises überdeckt. Schon lange hatte sie nicht mehr so stark gehofft. Wer nicht hoffte, konnte auch nicht enttäuscht werden. Und dennoch war Siantari nun enttäuscht, dass das Felsentor dem Beben standgehalten hatte.

Dafür hatte, zunächst unbemerkt von Siantari, etwas Anderes dem Einsturz Krahals nicht standgehalten. Mochten die Temm vor Jahrtausenden ihre eigene Magiequelle für das Felsentor genutzt haben, so war die unnatürliche Wolkenkuppel, die sie über die Gipfel des Kuolema gelegt hatten, durch uralte Drachenmagie gespeist worden. Denn auch die Drachen der Urzeiten hatten ein Interesse daran gehabt, dass die Eisdämonen im Fahlen Gebirge und ihre tiefe Schlucht abgeschirmt vom Rest der Welt blieben.

Nun, wo Krahal eingestürzt und die magischen Ströme der Drachen versiegt waren, lichteten sich die Wolken am Gipfel des Kuolema langsam im beißenden Wind. Der Himmel über der ewigen Eisfläche war wieder zu sehen. Zum ersten Mal in beinahe fünf Jahrhunderten erblickte Siantari wieder das Licht der Sonne. Ungewohnt prickelte es auf ihrer eisigen Haut.



\trennzeichen



Es gab auch noch einen anderen, der verstand, warum sich die Wolken von den Gipfeln des Kuolema verzogen. Haamun, der Hexer aus Andor. Er hatte bei seiner Überquerung des Gebirges die Magie der Temm gespürt, die das Felsentor für die Eisdämonen undurchdringlich machte. Und er hatte die Drachenmagie gespürt, die die Wolken auf den Gipfeln des Kuolema-Gebirges gehalten hatte.

Haamun trat vor seine Hütte und blickte zum Kuolema hinauf. Nun waren die schneebedeckten Gipfel der Berge sichtbar geworden. Der Himmel jenseits des Gebirges erglühte in rotem Licht. Da wusste Haamun, dass der letzte Drache besiegt worden war. Und wenn der Drache angegriffen hatte, so hieß das, dass der alte König tot war. Und dass es an der Zeit war, nach Hause zurückzukehren.

Haamun wurde von einem Hilferuf aus seinen Gedanken gerissen.

„Haamun!“, erklang die Stimme einer Minenarbeiterin, die aus einem Stollen zu rennen kam, „Haamun! Die Erde hat gebebt. Der Stollen hat nachgegeben. Seram hat eine Platzwunde am Kopf. Kannst du...“

Haamun hatte bereits eine Handvoll Heilkräuter aus der Steppe ergriffen und rannte los. In den Jahren, seit er hierhergekommen war, hatte er viele der seltenen und oft unbekannten Kräuter der Steppe gesammelt, Tränke gebraut und die verschiedensten Salben, Pulver und Gifte hergestellt. Wenn es einen Notfall gab, war er zur Stelle, um zu helfen. Selbst wenn den Geholfenen am Ende etwas Unvorhergesehenes geschah, wenn sie seit der Heilung grüne Haut hatten oder nur noch im Versmaß sprechen konnten... nun, solange Haamuns Hexerkunst ihr Leben rettete, war es das weitaus wert.

„Der Stollen ist nicht zusammengestürzt“, behauptete Seram trotzig, während Haamun seinen Kopf verarztete, „Im Gegenteil! Aus dem Felsen hat sich ein alter Stollen wieder befreit. Während ihr alle euch vor niederrieselnden Kieseln versteckt habt, habe ich gesehen, wie der Berg sich mir gegenüber öffnete! Einen uralten Gang mit Wänden voller Runen und ähnlicher Kritzeleien sah ich. Das Beben hat einen der uralten Gänge der Temm freigelegt!“

Haamun hielt inne. Er hatte schon seit langem gespürt, dass die Stollen unterhalb des Kuolema-Gebirges sein Weg zurück nach Andor sein könnten. Nun, mit dem Tod des alten Königs und des Drachen, begann eine neue Ära für Andor. Eine, an der Meres wieder teilhaben konnte. Und ausgerechnet jetzt wurde ein uralter Temm-Stollen freigelegt? Alle Spielsteine fielen an ihre Stelle. Dieser Gang würde sie nach Andor führen.

Es war an der Zeit. Nun musste er rasch handeln. Mitreisende auftreiben für die Rückkehr in seine Heimat. Haamun selbst, seine Tränke und seine Geschichten waren beliebt bei einigen Bauern nahe des Gebirges, insbesondere den jüngeren. Vielleicht würden sich einige davon ihm anschließen. Und in den nahe gelegenen Dörfern konnte er sicherlich einige Händler, Baumeister und Kartographen zusammentreiben, deren Neugier stärker war als die Furcht vor den Geistern der Berge. Mit einer ganzen Gesellschaft würde er in Andor aufkreuzen und in den unruhigen Zeiten nach dem Tod von König und Drache aushelfen. Sie würden ihm verzeihen und den längst vergangenen Zwist beilegen. Er konnte wieder Meres sein.



\trennzeichen



Aćh war gerade unterwegs in der Metropole Agarb, als sie ein Falke mit einer Nachricht von Barz erreichte.\bigskip



„Hochverehrte Hüterin Aćh,



Ein Beben hat am Ende eines Stollens der Bergarbeiter eine Verbindung zu einem uralten Gang der Temm freigelegt. Haamun, der Hexer, ist sich absolut sicher, dass er uns nach Andor führen kann. Er trommelt eine Gruppe Reisender zusammen, die sich für das Königreich auf der anderen Seite des Kuolema interessieren. Und der Gang scheint gerade groß genug für Sabri zu sein. Ich habe vor, mich ihm anzuschließen. Bist du dabei? Wir brechen in wenigen Tagen auf.



Beste Grüße

Barz“\bigskip



Wie weit Barz‘ Sprachkenntnisse doch in den letzten Jahren gewachsen waren!

Aćh schluckte schwer und ließ ihre Optionen Revue passieren. Endlich war es soweit. Seit dem Gespräch mit dem Hüter der Zeit damals hatte Barz gewusst, dass es einen Weg für ihn nach Andor und zurück zu seiner Familie geben würde. Und in den Jahren seit damals, in denen sie und Barz sich näher gekommen waren, war mehrmals die Frage aufgekommen, ob sie ihn begleiten würde.

Sie mochte ihre Arbeit am Nestbaum der Takuri, und sie mochte ihre Freunde und Bekannte dort. Doch sie mochte auch Barz, und sie mochte Abenteuer. Sturmgeister und Eisdämonen, auf so etwas stieß man nicht, wenn man sein Leben lang am selben Baum verbrachte. Andor, das fremde Land, konnte bestimmt von diplomatischem Kontakt mit Tulgor profitieren, und als Tochter einer Diplomatin war Aćh für eine solche Rolle vermutlich besser geeignet als ein Hexer und seine zusammengewürfelte Reisegruppe.

Noch während Aćh ihre Gedanken sortierte, klopfte es an ihrer Tür. Als sie sie öffnete, blickte sie überrascht in das Gesicht ihrer Mutter... und in die verschmitzten Augen des Hüters der Zeit, der auf ihrem Rücken saß. Seinen Turban hatte er gegen eine braune Kutte eingetauscht.

„Na, schon reisefertig gemacht?“, krächzte der Temm grinsend, „Ich schließe mich dir gleich an, wenn das für dich in Ordnung ist. Schließlich will ich auch mit nach Andor!“



\trennzeichen



Gemeinsam hatte Aćh mit dem Takuri das Land durchquert und vor wenigen Tagen das Gebirge erreicht. Die Sonne war hinter den schneebedeckten Bergspitzen untergegangen, und ein roter Mond stand am Himmel. Nun sah der Takuri Aćh noch einmal an, und in einem kurzen Augenblick rotglühenden Feuers beendete er sein Leben, so wie er es schon so oft beendet hatte. Als kleines Küken watschelte er in Aćhs Arme und kuschelte sich an sie.

Der Hüter der Zeit betrachtete Turr mit zusammengekniffenen Augen:

„Ist alles in Ordnung mit dem Takuri? Er ist nun schon zum zweiten Mal in zwei Wochen verglüht.“

Aćh schüttelte ihren Kopf.

„Nein, irgendetwas ist besonders an ihm. Seit seinem Verschwinden und seiner Zeit im ewigen Eis altert er viel zu schnell. Manchmal durchlebt er seinen Zyklus in wenigen Tagen!“

„Vielleicht finden wir ja in Andor etwas, das ihm helfen kann“, meinte Barz optimistisch, und setzte sich neben die beiden ans Lagerfeuer, „Und daran zu leiden scheint er auch nicht.“

Mit einem nachdenklichen Blick setzte Barz nach: „Ist das eigentlich nach jedem Zyklus immer noch derselbe Takuri?“

Aćh grinste. Sie hatte sich eine solche Frage schon oft gestellt. War ein Takuri nach seiner Regeneration noch derselbe Takuri? Wie sollte man sie nur beantworten? Der neue Takuri würde nie identisch aussehen wie der alte, aber dennoch würde der neue Takuri dem alten erheblich ähnlicher sehen als allen anderen. Vielleicht mit einer anderen Federanordnung oder sonstigen Kleinigkeiten. Auch wenn das Küken noch auf den alten Namen reagieren würde, stellte sich die Frage, ob esdieselben Erinnerungen wie der frühere Takuri trug. Aćh hatte oft genug mit einem Töpfchentraining von vorne anfangen müssen. Gewisse Eigenheiten, bestimmte Verhaltensmuster, blieben dennoch nach jeder Regeneration, soweit sie sagen konnte. Wobei ihre Einschätzung natürlich vielen Vorurteilen unterlag. Wie wahrheitsgetreu sie war, war schwer einzuschätzen.

Aćh beschloss, Barz eine Gegenfrage zu stellen: „Was würdest du sagen: Wenn du am Morgen aufwachst, bist du dann derselbe Mensch, der eingeschlafen ist?”

„Natürlich”, meint Barz, „Wenn wir uns nach jedem Schlaf als andere Person sehen würden, gäbe es ja kaum einen Grund, uns groß um die nächsten Tage zu sorgen. Wir würden doch einfach unseren einzelnen Tag mit kurzweiligen Genüssen verschwenden und die Menschheit hätte sich gar nie erst zu einer Gesellschaft entwickelt.”

Das war nicht die Sorte Antwort, die Aćh erwartet hatte.

„Pfff. Da kann ich ja nicht einmal eine Analogie zur Regeneration des Takuri ziehen. Außerdem verzichten wir doch für unsere Nachkommen auf gewisse Genüsse, könnten wir das dann nicht auch hier tun? Und dass etwas ohne den Glauben an etwas nicht entstanden wäre, sagt ohnehin nichts darüber aus, wie wahr dieses etwas wirklich ist, oder etwa nicht?”

Barz legte den Kopf schief: „Da magst du Recht haben. Ich weiß, dass du diesen Takuri hier nach jeder Regeneration immer noch Turr nennst. Macht ihr das normalerweise immer so, oder gebt ihr den Älteren jeweils neue Namen, wenn sie sich regenerieren?”

„Wir sehen von Fall zu Fall, ob wir einen neuen Namen verleihen wollen. Oftmals ist es praktischer, einfach dieselben Namen zu verwenden. Manchmal ist ihre Persönlichkeit nach einem Zyklus aber einfach zu anders, und dann wählen manche Hüter neue Namen. Dann gibt es oft eine Namenszeremonie, wie wenn einer von uns den seinen wechselt.“

„Tulgori wechseln ihre Namen?”

„Ihr Barbaren etwa nicht?“

„Natürlich nicht!“

„Hast du denn in all deiner Zeit hier keine Namenszeremonie mitgekriegt? Heißt du als Kind etwa gleich wie als Jugendlicher? Wie als Volljähriger? Wie als alter Erwachsener?”

„Natürlich heiße ich immer gleich.”

„Aber ein Kind kann doch noch keinen Namen wählen. Das Kind kennt die Bedeutung ja gar nicht.”

„Das muss es auch nicht. Die Eltern wählen ja für das Kind.“

„Und was, wenn die Eltern sich nicht einig sind?“

„Dann diskutieren sie das halt aus.“

„Und was, wenn das Kind den Namen später nicht mag?“

„Jeder, den ich kenne, mag seinen Namen. Irgendwie.“

„Was, wenn nicht? Ich mochte meinen Geburtsnamen nicht.“

„Du wusstest ja auch, dass du ihn ändern kannst. Wir gewöhnten uns dran.“

„Na eben.“

„Dann weiß man ja nicht mal, von wem man redet, wenn Leute später immer anders heißen.“

„So oft ändert sich der Name ja auch wieder nicht. Und man gewöhnt sich daran.“

Hinter Barz grunzte es, als seine Echse Sabri ihn endlich erreichte (sie hatte sich in der letzten Stunde unglaublich gemächlich auf ihn zubewegt) und sich ächzend auf den Boden sinken ließ. In den letzten Jahren war sie beachtlich gewachsen, sodass sie ihm im Stehen nun schon bis zur Schulter reichte. Sie war über und über beladen mit verschiedenen Pulversäcklein. Barz‘ Vorräte mochten beim Kampf gegen den Sturmgeist erheblich reduziert worden sein, doch hatte er in der Steppe Tulgors und in zahlreichen Gesprächen mit dem Hexer Haamun seine Pulversammlung teils aufstocken (Mondbeeren für sein Bannpulver schienen beispielsweise buchstäblich überall zu wachsen) und teils gar alternative Pulvermischungen zusammenstellen können.

Haamun und sein Trupp waren gleich zwei Tage nach dem großen Beben, dem Verschwinden der Wolken über den Gipfeln und dem Freilegen des Temm-Stollens in der Mine verschwunden, um die Lage mit eigenen Augen auszukundschaften. Schon am ersten Abend waren sie mit glänzenden Augen und einer Handvoll wertvoller Edelsteine in seiner Tasche wieder beim Stolleneingang aufgetaucht und berichteten von einer „lohnenswerten Sackgasse“.

Nun, eine ganze Woche nach dem großen Beben, war die Kompanie schon zum dritten Mal an den Startpunkt zurückgekehrt. Ihre gute Laune hatte sich gelegt und Haamun verbrachte viel Zeit damit, mit einer Art Wünschelrute in seinen Händen im Eingangsbereich verschiedener Stollenabzweigungen umherzulaufen. Die unterirdischen Gänge der Temm hatten sich als wahres Labyrinth herausgestellt, welches nur mit der nötigen Magiekunde navigiert werden konnte. Und Haamun rang mit seinem Stolz.

Eforas, einer seiner engsten Vertrauten, unterhielt sich soeben etwas abseits des Lagerfeuers mit ihm. Man konnte von außen nicht viel vernehmen, doch Aćh glaubte eindeutig, ein „Du sagtest, du würdest den Weg durch den Berg kennen!“ von Eforas zu vernehmen, sowie eine zischende Antwort von Haamun.

Ein Raunen ging durch die Dutzend Händler, Baumeister und Kartographen, die ebenfalls in der Nähe saßen.

Aćh wisperte zum Hüter der Zeit: „Ich weiß, dass diese Temm-Gänge viele Verzweigungen besitzen, aber du bist ein Temm und kannst noch dazu die Zukunft sehen! Kannst nicht einfach du uns durch die Gänge führen?“

„Wer kann schon wissen, was genau ich alles tun kann?“, meinte der Hüter mit einem verlegenen Grinsen, „Aber das würde Haamun wohl nicht gefallen. Bald schon wird er von selbst den richtigen Einfall haben. Und ganz abgesehen davon mag ich die Gesellschaft hier.“

Aćh und Barz blickten sich an und grinsten. In Barz‘ Blick war aber auch eine lange unterdrückte Sehnsucht zu erkennen. Gedankenverloren griff er an seinen Hals und spielte mit einer goldenen Kette, die Aćh ihm vor einem Jahr geschenkt hatte. Barz‘ Ringkette war, soweit sie wussten, noch immer vor dem Felsentor der Eisdämonen eingeschneit.



\trennzeichen



„Was... was ist das?“, stammelte Haamun leise.

„Es ist wunderschön“, murmelte Aćh.

Haamun hatte endlich – so hofften alle – die richtigen Abzweigungen gefunden. Schon seit einigen Tagen waren sie nun die Stollen unter dem Berg entlanggewandert, ohne sich in Sackgassen zu verirren. So langsam wünschten sich alle Reisenden, wieder das Licht der Sonne zu sehen und den Wind in den Haaren zu spüren. Doch der Anblick, der sich ihnen hier unter der Erde bot, war beinahe ein Ausgleich dafür.

Sie hatten eine Höhle unter der Erde erreicht, aus deren Wände dutzende Edel- und Mera-Steine ragten. Ein leises Summen erfüllte den Raum, und wie überall sonst waren auch hier zahlreiche Runen am Boden zu sehen. Turr entfaltete seine Flügel und erhob sich zum ersten Mal in einigen Tagen zu einem kurzen Rundflug in der geräumigen Höhle. Er stieß einen Freudenschrei aus. Seine Schwanzfedern glommen glücklich.

Barz klopfte der dem Trupp langsam nachtrottenden Sabri auf den Rücken. Sie folgte dem Angebot bereitwillig und ließ sich zu Boden plumpsen. Schon bald ertönte ein Schnarchen aus ihrem breiten Mund.

Der Hüter der Zeit kniff seine Augen zusammen und runzelte seine Stirn. Nach einiger Zeit öffnete er sie wieder und sprach:

„Tut mir leid, werte Reisende. Ich habe keine Ahnung, um was für eine magische Höhle es sich hier handelt. Ich werde dies mein ganzes Leben lang nicht in Erfahrung bringen. Und ihr demnach wohl ebenfalls nicht.“

Murmeln war unter den Reisenden zu hören. Dann sprach Haamun leise und monton:

„Es wäre schon fast schade, dieses schöne Bild zu zerstören und die Mera-Steine mitzunehmen.“

Eforas nickte: „Fast.“

Energisch trat er nach vorne an die Wand und untersuchte einen Mera-Stein, der darin verankert war. Wie üblich sah es aus, als wäre der farbige Stein mit dem umliegenden Felsen verschmolzen.

„Sitzt fest wie angeleimt.“

Eforas zückte einen Pickel und ließ ihn auf den Felsen rund um den Stein niederfahren. Der Pickel zerbrach mit einem hellen Klingen. Es erinnerte Aćh an daran, wie die Werkzeuge der Minenarbeiter auch am Nestbaum der Takuri zerschellten.

„Hier wollte jemand nicht, dass diese Schätze entwendet werden“, meinte Eforas schicksalsergeben.

„Schaut mal her!“, rief Barz da. Er hatte in den Rucksack auf Sabris Rücken gelangt, den die Reisenden mit allerlei Schätzen und Gaben beladen hatten, und ein rautenförmiges Felsstück herausgezogen. Aus dem Felsstück ragte ein schwach grünlich leuchtender Mera-Stein. In dieser Form wurden die Steine üblicherweise weiter versandt, um weiter Inlands komplett aus dem Felsen geschält zu werden. Ein komplexer Prozess, der den Minenarbeitern abgenommen wurde.

Barz hob den grünen Mera-Stein in die Höhe und brachte ihn in die Nähe eines in die Höhlenwand eingelassenen grünen Mera-Steins. Je näher Barz trat, desto heller leuchteten die beiden magischen Steine, und ein immer lauteres Summen war zu vernehmen. Während Barz die Distanz zwischen den Mera-Steinen verringerte, wurde das Summen zudem immer dumpfer und tiefer, bis die Reisenden es nicht mehr hören konnten, doch die Vibrationen immer noch in ihrem Körper zu fühlen glaubten.

Dann machte es „klick“, als die beiden Mera-Steine einander berührten. Das helle Leuchten erlosch – und der Mera-Stein aus der Wand rutschte herunter und kullerte über den Boden, als hätte ihn nie etwas in der Wand gehalten.

Da waren die Reisenden aber außer Rand und Band. Innert Minuten hatten sie alle Mera-Steine mithilfe anderer Mera-Steine aus der Wand gelöst und eingepackt. Nachdem jemand feststellte, dass sich die Edelsteine ebenfalls so lösen ließen, dauerte es sogar noch kürzer, bis auch diese allesamt ordentlich in einem Sack versammelt waren. Außer die ein, zwei Steinchen, die vielleicht jemand lieber in seine eigene Hosentasche hatte wandern lassen. Sabri wurde eine weitere Tragetasche voller Mera-Steine übergehängt. Die Echse ächzte, wehrte sich aber nicht.

In vielen Augen blitzte Begeisterung über den unerwarteten Fund. Der Hüter der Zeit blickte hingegen nur nachdenklich drein. Und auch Haamun fiel nicht in Freudenschreie ein. Ausgerechnet Meres machte sich Gedanken darüber, ob das Lösen der Mera-Steine ungewollte Konsequenzen haben könnte.



\trennzeichen



Siantari tigerte hinter dem Felsentor über das ewige Eis. Hin und her, hin und her. Ihre Augen waren starr auf den Himmel gerichtet. Die Wolken hatten sich verzogen, und die Sterne zogen ihren Blick magisch an. Die klirrende Kälte, die Siantari unerbittlich entgegenschlug, konnte ihr nichts anhaben. Sie war die frostigen Tage und Nächte im Kuolema-Gebirge gewohnt.

So stand sie da und blickte in den weiten Kosmos, weiß auf Blau. Ein Kosmos, in dem sie und ihre Eisdämonen doch nur ein so kleiner, unwichtiger Teil waren. Noch. Denn eines Tages würde das ewige Eis über alles unter dem Himmelszelt gewachsen sein. Und Siantari würde dabei geholfen haben.

Sie verspürte Hoffnung. Ja, die Sterne erfüllten sie mit Hoffnung.

Die uralte Drachenmagie, die die Gipfel des Kuolema stets in Wolken verhüllt hatte, war verklungen. Alles verfiel mit der Zeit. Auch uralte Magien. Und so konnte es nicht mehr unendlich lange dauern, bis auch ihr Felsentor geöffnet wurde und sie oder einer ihrer Nachfolger das ewige Eis verlassen konnte.

Etwas regte sich in ihrem Hinterkopf. Eine Wahrnehmung, klein und fein, und doch so unglaublich relevant.

Sie fühlte, dass der Fremde zurückgekehrt war. Der Hexer mit den grünen Flammen, der das ewige Eis vor einigen Jahren, die für Siantari kaum mehr als ein Wimpernschlag her waren, vom Osten her überquert hatte. Doch überquerte er den Kuolema nicht. Er ging unter den Bergen hindurch, und er war nicht allein. Und sie trugen Mera-Steine bei sich, diese Narren!

Siantari spürte, wie sich etwas grundlegend änderte. Damals, als sie zur Eisdämonin geworden war, hatte sie auf einmal eine Erkenntnis ergriffen. Dass sie nie wieder durchs Felsentor gehen konnte. Dass sie nun eine Gefangene hinter dem Felsentor war, dazu verdammt, 500 Jahre lang das ewige Eis zu nähren, um am Ende einen neuen Eisdämon zu erschaffen und dann für immer im Eis zu versinken. Diese Erkenntnis, magisch auferlegt durch das Felsentor und alle anderen uralten Vorrichtungen der Temm, die die Eisdämonen von allen Seiten in diesen Schluchten einsperrten... noch vor wenigen Augenblicken war diese Erkenntnis Siantari noch bewusst gewesen, wie eine leise Last, die ihr seit Jahrhunderten auf den Schultern lag. Nun war sie einfach nicht mehr da. Sie fühlte sich so unbeschwert. Der Bann war verschwunden.

Misstrauisch schritt Siantari auf das Felsentor zu. Konnte es sein, dass das Tor geöffnet worden war? Es sah noch genau gleich aus wie zuvor. Doch fühlte es sich ganz anders an.

Triumphierend glitt Siantari nach vorne und marschierte durch das Felsentor. Keine uralte Magie hinderte sie daran. Die Ströme aus dem Innern des Bergs waren versiegt.

Der Fremde und seine Begleiter mussten die Quelle des Felsentors gefunden haben. Und ihre Mera-Steine hatten den Mechanismus ausgehebelt. Sie hatten das Felsentor geöffnet. Siantaris Felsentor war offen. Sie war frei!

Ein kurzes Lächeln huschte über Siantaris Lippen.

Und es kam noch besser. Die Reisenden tief unter dem Gebirge trugen weiterhin Mera-Steine mit sich. Siantari konnte die Steine und ihre Lage spüren, so wie sie den gesamten Berg spüren konnte. Schließlich war das ihr Berg. Das waren ihre Steine. Und so würde sie stets wissen, wohin diese Reisenden reisten. Sie konnte ihnen folgen. Sie würden ihr den Weg durch den Berg zeigen. Der Weg in ein fremdes Land. Sollte sie ihnen folgen?

Siantari blickte sich um. Links oder rechts? Wollte sie zurück nach Tulgor, in das Land ihrer Kindheit, oder lieber in den Osten, wohin diese Narren reisten und woher dieser mysteriöse Hexer gekommen war?

Leise Wehmut übermannte Siantari, als sie an ihre Heimat dachte. Ein überaus menschliches Gefühl, wie sie überrascht feststellte. Wollte sie dorthin zurückkehren? Sehen, was aus der Hütte ihrer Familie geworden war?

Nein, sprach sie zu sich. Tulgor bedeutete ihr nichts. Zeit, neue Wege zu beschreiten. Zeit, sich der Zukunft zuzuwenden. Zeit, in das Land im Osten aufzubrechen.

Siantari erhob sich und blickte ein letztes Mal über das ewige Eis des Kuolema. Hier und dort glaubte sie, ferne Silhouetten anderer Eisdämonen ausmachen zu können. Einsame, verwirrte Gestalten.

Siantari war es als erster Dämonin des ewigen Eises gelungen, Eiskristallketten zu fertigen, die ihre Gabe und ihre Mission an andere weitergeben können. Hin und wieder hatte sie aus verirrten Bergsteigern Durs und Doras geschaffen. Sie hatte sie einsetzen wollen, um aus dem Felsentor auszubrechen. Doch hatte die lange Zeit im ewigen Eis den menschlicheren Geistern der Durs und Doras nicht gutgetan. Und nun brauchte Siantari brauchte sie nicht mehr. Das ewige Eis konnte sie auch allein bis in die umliegenden Länder ausbreiten.

„Meine Kinder“, erhob sie ihre kalte Stimme, und das Echo ihrer Worte schallte über die riesige Eisfläche, „Ihr seid nun endlich frei, und frei sollt ihr nun sein. Ich entlasse euch aus meinen Diensten. Geht, wohin euch eure Schicksale ziehen. Verbreitet das ewige Eis. Wenn es so sein soll, werden wir uns wiedersehen.“

Keine Reaktion kam aus dem ewigen Eis. Siantari wusste, dass ihre Worte angekommen waren. Doch wie viele der Eisdämonen hier besaßen überhaupt noch genug Willenskraft, um sich zu irgendeiner Handlung aufzurappeln? Sie war allein. Aber das war in Ordnung. Sie war schon so lange allein gewesen. Sie wandte sich nach Osten und verließ das ewige Eis und das Fahle Gebirge. Zeit, alles Land in eine Eiswüste zu verwandeln.

\begin{commentbox}
    Weiter geht es in \hyperref[2022-runen-im-schnee]{Runen im Schnee (2022)}.
\end{commentbox}